\section{Running Example: Economic-Policy State}\label{sec:min-manifesto-state}

The Benoit/Manifesto rubric scores six policy dimensions; we keep one, the
economic left--right score. For a manifesto span $x$, let
$S_{\mathrm{econ}}(x)$ denote the task-relative state containing the
evidence a trusted coder would use to place the span on that scale:
commitments to public services, taxation, welfare
expansion or retrenchment, market intervention, privatization, fiscal
restraint, and the qualifications that determine how those statements should
be read. It is the economic analogue of a sufficient statistic for the
rubric: it may discard slogans and style, but it must keep anything that can
move the economic score. The paper proceeds as if this state exists---that
is a modelling assumption about the measurement problem, not a result---and
the audit returns an estimate of how badly the assumption is violated on
the realized tree.

The corresponding oracle is a readout of that state:
\[
  \fstar_{\mathrm{econ}}(x)
  = r_{\mathrm{econ}}\!\left( S_{\mathrm{econ}}(x) \right)
  \in \{1, \ldots, 7\},
\]
where the seven-point scale is the Benoit/Manifesto expert-survey scale
\citep{BenoitEtAl2025}. The Manifesto/Benoit benchmark itself is a
\emph{root-observed} corpus evaluation: a sample of party platforms, each
already carrying a full-document expert target for the economic dimension;
the tree root is the predicted score for that platform, scored directly
against the expert target. In the results reported below,
$r_{\mathrm{econ}}$ is approximated in practice by a rubric scorer applied
to the root summary---a language model acting under a fixed scoring prompt
that returns a seven-point integer, not a free-form judgment. The target,
however, is always the same: the expert score for the full document.

The Manifesto Project also ships a span-level oracle: expert
quasi-sentence codings for $2{,}157$ platforms, about $2.27$M coded spans
in total. Aggregating coded quasi-sentences inside a leaf or merge span,
and reading out the corresponding economic score, instantiates oracle
values at paragraph and section granularity---the ingredients the
node-level audit in Section~\ref{sec:min-audit} needs. These labels
are the planned node-level route for this setting; they are not the
evidence behind the root-observed headline of
Section~\ref{sec:min-manifesto-results}, which stands on its own as a
full-document evaluation.

A produced summary is useful only insofar as it acts like a compact
economic state. A leaf summary is faithful if $\fstar_{\mathrm{econ}}$
gives the same reading on it as it does on the raw span---the summary
can drop slogans, but it must keep the evidence the rubric would score.
A merge of two sibling summaries is faithful if it preserves the economic
state of their union, including cross-section tensions that neither
child can resolve alone: a platform that is environmentally
interventionist in one section but fiscally restrained in another has
its economic reading at the boundary between those sections, not in
either one. And a second summarization pass---feeding a stored summary
back through the summarizer---should not shift the economic readout.
Section~\ref{sec:min-framework} writes these three requirements as local
laws on strings, keeping $\fstar_{\mathrm{econ}}$ in place as the task
oracle throughout.

The economic state is \emph{latent}: no dataset labels $S_{\mathrm{econ}}(x)$
itself, and no isolated unit test can certify that a given summary preserved
it under $\fstar_{\mathrm{econ}}$. To separate ``the framework is sound''
from ``the summarizer did its job on a latent target,'' the paper validates
the mechanism on two tasks whose states are already known and whose merge
rules can be written by hand. Markov changepoint counting anchors the
ordered-state case: the summary must carry boundary information or the count
is systematically off. HyperLogLog anchors the unordered mergeable-sketch
case: the scalar distinct-count answer is not mergeable, but a bounded register
state is. In both anchors, a known oracle sits on the other side of the learned
mechanism, so a failure has to be the framework's, not the summarizer's.

The formal artifact follows the same boundary. It mechanizes the state-level
merge laws used by the anchors, including ordered list-homomorphism schedule
invariance, HLL's max-register merge, Count-Min's non-idempotent additive merge,
Misra--Gries/SpaceSaving bookkeeping, finite VC trace union bounds, and the
classical sketch/readout interface. The heavier literature facts---HLL
asymptotic constants, the compact randomized range/quantile constructions,
geometric $\varepsilon$-kernel size proofs, and MUD communication lower
bounds---are recorded as typed theorem obligations, not as informal
assumptions. The standalone Lean
entry point for this layer is
\texttt{FormalProbability.ML.MergeableSummaries.Literature}; the C-TreePO
formalization imports that surface and re-exports the paper-facing bridge
names. For the HLL running example, the bridge also checks the corrected
precision arithmetic $1.04/\sqrt{2^{14}}=13/1600<1\%$.

A reader who prefers plays to platforms can follow the same contrast
through \emph{Hamlet}. \emph{How many shifts of emotional register occur
across the play} is the Markov analogue. Each scene sits in one of a small
set of registers---Mourning, Foreboding, Political, Intrigue, Madness,
Action, Comic---and a shift at a scene or act boundary is invisible to either
act alone unless each act-summary records its first and last register, so the
merge rule can add a boundary indicator when they disagree. \emph{How many
distinct character co-appearance pairs ever share a scene} is the HLL
analogue: each scene emits the set of unordered speaker pairs on stage, each
pair is one item in the stream, each act merges by the union of those
pair-sets, and HLL's register array with
pointwise-max merge approximates union cardinality without storing the
union-set. (The simpler ``how many distinct speaking roles appear''---about
twenty-four---is the enumerable warm-up; HLL earns its place on a corpus of
plays, where the union grows while the register array stays fixed-size.)
These examples are deliberately narrower than literary or political
interpretation; they are useful because a competent student can write the
merge rule down. The manifesto question is harder precisely because the
\emph{state itself}---what an economic-policy coder would keep, and at what
granularity---has to be learned from examples and then audited.
Section~\ref{sec:min-validation} returns to these tractable cases to show
the mechanism preserves the oracle value the hand-written merge would
preserve; Section~\ref{sec:min-manifesto-results} reports the learned case on
the Manifesto/Benoit benchmark.
